\chapter{Recalls of Linear Algebra}

\begin{definition}[eigenvalues and eigenvectors]
Let $V$ be a vector space over the field $\KK$ (with $\KK = \RR$ or $\KK=\CC$)
and $A\colon V\to V$ a linear operator.
We say that $\lambda\in \KK$ is an \emph{eigenvalue}%
\mymargin{eigenvalue}\index{eigenvalue} of $A$
if there exists $v\in V$, $v\neq 0$ such that
\[
  Av = \lambda v.
\]
In this case, $v$ is said to be an \emph{eigenvector}%
\mymargin{eigenvector}\index{eigenvector} of $A$ corresponding
to the eigenvalue $\lambda$.

We denote by $A-\lambda$ the linear operator $A-\lambda I$
where $I\colon V\to V$ is the identity.
The vector space
\[
  \ker (A-\lambda I)
\]
is called the \emph{eigenspace}%
\mymargin{eigenspace}\index{eigenspace} corresponding to the eigenvalue $\lambda$.
The eigenspace consists of the zero vector and all eigenvectors
corresponding to the eigenvalue $\lambda$.

We say that $v$ is a generalized eigenvector of degree $m$ ($m\ge 1$ integer)
corresponding to the eigenvalue $\lambda$
if
\[
 (D-\lambda)^m v = 0
 \qquad \text{but} \qquad
  (D-\lambda)^{m-1} v \neq 0.
\]

Generalized eigenvectors of degree $1$ are precisely the
eigenvectors.
\end{definition}

\begin{proposition}[properties of eigenvectors]
Let $A\colon V \to V$ be a linear operator.
We denote by
\[
  W_\lambda^m = \ker (A-\lambda)^m \setminus \ker (A-\lambda)^{m-1}
\]
the set (not a vector space!) of all generalized eigenvectors
of $A$ of degree $m$ corresponding to the eigenvalue $\lambda$.
Then if $\lambda \neq \mu$ we have
\begin{enumerate}
\item eigenvectors corresponding to distinct eigenvalues are distinct:
\[
W_\lambda^1 \cap W_\mu^1 = \emptyset;
\]
\item the operators $A-\lambda$ and $A-\mu$ commute:
\[
  (A-\lambda)(A-\mu)v = (A-\mu)(A-\lambda);
\]
\item
the operator $A-\mu$ leaves the sets $W_\lambda^m$ invariant:
\[
  v\in W_\lambda^n
  \implies
  (D-\mu)v \in W_\lambda^n.
\]
\end{enumerate}
\end{proposition}
%
\begin{proof}
\begin{enumerate}
\item
If there existed $v\in W_\lambda^1 \cap W_\mu^1$, we would have
\[
(\lambda - \mu) v = \lambda v - \mu v = Av - Av = 0.
\]
But this is impossible if $v\neq 0$ and $\mu\neq \lambda$.

\item
We have
\[
  (A-\lambda)(A-\mu)
  = A(A-\mu I) - \lambda (A-\mu I)
  = A^2 - \mu A -\lambda A + \lambda \mu I
  = A^2 - (\mu+\lambda) A + \lambda\mu I.
\]
Since the right-hand side of the equality is invariant if
we swap $\lambda$ and $\mu$, the left-hand side must also be invariant.

\item
If $(A-\lambda)^m v = 0$, we have
\[
(A-\lambda)^m (A-\mu) v = (A-\mu)(A-\lambda^m) v = 0.
\]
We want to show that if $(A-\lambda)^{m-1} v \neq 0$
then also $(A-\lambda)^{m-1} v \neq 0$.

\end{enumerate}
\end{proof}